%!TEX root = hems_proj.tex
%%%%%%%%%%%%%%%%%%%%%%%%%%%%%%%%%%%%%%%%%%%%%%%%%%%%%%%%%%%%%%%%%%%%%%%
\chapter{Kísérleti Eredmények}\label{ch:GYAK}
%%%%%%%%%%%%%%%%%%%%%%%%%%%%%%%%%%%%%%%%%%%%%%%%%%%%%%%%%%%%%%%%%%%%%%%

\begin{osszefoglal}
Ez a fejezet két szintet fed le: a klasszikus háromszcenáriós (tél/nyár/borús) eredményeket, valamint a 2×2-es kísérleti mátrix benchmark protokollt. A cél a korábbi eredmények megőrzése mellett a teljes pipeline tudományos értelmezése, beleértve az alkalmazott paramétereket is.
\end{osszefoglal}

\section{Kísérleti szcenáriók}

A kísérleteket három, markánsan eltérő szcenárióban futtattuk le 20-20 alkalommal, hogy a rendszer robusztusságát vizsgáljuk. A téli esetet az alacsony PV termelés és a magasabb alapfogyasztás jellemzi, ami nagy hálózati függést eredményez. A nyári időszakban ezzel szemben magas, csúcsban akár 6 kW-os PV termeléssel és alacsonyabb fogyasztással számolhatunk, ami jelentős visszatáplálási potenciált rejt. Harmadikként egy viharos/borús napot szimuláltunk, ahol az ingadozó PV termelés és a magas zaj miatt az akkumulátor dinamikus kezelése válik szükségessé, ezáltal stressztesztnek vetve alá az algoritmusokat. A reprodukálhatóság érdekében minden futás determinisztikus bázis-seeddel (42) indul, amelyet a 20 ismétléshez futásonként eltoltunk, a nyers eredményeket és a futási idők listáit pedig CSV/JSON formában mentettük.

\section{Paraméterek beállítása}

Az algoritmusok összehasonlíthatóságát egységes paraméterezési kerettel biztosítottuk. Minden algoritmus azonos populációméreten és iterációszámon futott: $N = 50$ (illetve GA esetén $N = 100$ a változat-összehasonlításhoz) és $G = 150$ iteráció. A döntési vektor dimenziója $d = 26$ volt minden esetben.

A GA-nál a mutációs valószínűség $p_m = 0.1$, a keresztezés aritmetikai átlagolással történt, a szülőválasztás tournament szelekció (3 versenyző) alapján. Az elitizmus mértéke 2 egyed generációnként. A PSO-nál a tehetetlenségi súly $w = 0.6$, a kognitív gyorsulási konstans $c_1 = 1.8$, a szociális gyorsulási konstans $c_2 = 1.8$. A GWO-nál az $a$ paraméter lineárisan csökkent $2$-ről $0$-ra az iterációk során, más hangolható paraméter nem volt szükséges. A hibrid algoritmusnál a PSO-fázis az iterációk első $60\%$-át tette ki ($\lfloor 0.6 \cdot G \rfloor = 90$ iteráció), a GA-fázis a maradék $40\%$-ot ($60$ iteráció). A seed injection során a populáció fele a PSO legjobb megoldásának Gauss-perturbációiból ($\sigma = 0.5$) állt.

Az akkumulátor modelljében a kapacitás $C_{batt} = 10$ kWh, a SOC-határok $[0.2, 0.9]$, a maximális töltési/kisütési teljesítmény $P_{max} = 3$ kW. A büntetési együtthatók $\lambda_{soc} = 5000$ és $\lambda_{end} = 2000$ voltak. A visszatáplálási ár az aktuális vételezési ár $30\%$-a.

\section{Téli szcenárió eredményei}

Az alapvonal (akku és időzítés nélkül) télen átlagosan \textbf{42.09 RON}/nap.
\begin{table}[H]
\centering
\caption{Téli szcenárió statisztikái (20 futás, 50/100-as populáció, 150 iteráció, globális seed: 42)}
\label{tab:winter_stats}
\begin{tabular}{lccc}
\toprule
\textbf{Algoritmus} & \textbf{Átlag költség} & \textbf{Legjobb futás} & \textbf{Átlagos idő (s)} \\
\midrule
GA (P=50)  & 37.05 & 34.92 & 0.31 \\
GA (P=100) & 37.25 & 35.06 & 0.62 \\
PSO        & 36.87 & 34.56 & 0.21 \\
\textbf{GWO}    & \textbf{35.40} & 33.91 & 0.44 \\
Hybrid     & 36.28 & \textbf{34.42} & 0.24 \\
\bottomrule
\end{tabular}
\end{table}

\section{Nyári szcenárió eredményei}

Az alapvonal (akku és időzítés nélkül) nyáron átlagosan \textbf{12.48 RON}/nap.
\begin{table}[H]
\centering
\caption{Nyári szcenárió statisztikái (20 futás, 50/100-as populáció, 150 iteráció, globális seed: 42)}
\label{tab:summer_stats}
\begin{tabular}{lccc}
\toprule
\textbf{Algoritmus} & \textbf{Átlag költség} & \textbf{Legjobb futás} & \textbf{Átlagos idő (s)} \\
\midrule
GA (P=50)           & 4.08 & 3.74 & 0.31 \\
\textbf{GA (P=100)} & \textbf{3.93} & \textbf{3.64} & 0.62 \\
PSO                 & 6.67 & 5.20 & 0.20 \\
GWO                 & 5.17 & 4.21 & 0.43 \\
Hybrid              & 5.88 & 4.93 & 0.25 \\
\bottomrule
\end{tabular}
\end{table}

\section{Viharos/Borús szcenárió eredményei}

\begin{table}[H]
\centering
\caption{Viharos/Borús Nap szcenárió statisztikái (20 futás, 50/100-as populáció, 150 iteráció, globális seed: 42)}
\label{tab:cloudy_stats}
\begin{tabular}{lccc}
\toprule
\textbf{Algoritmus} & \textbf{Átlag költség} & \textbf{Legjobb futás} & \textbf{Átlagos idő (s)} \\
\midrule
GA (P=50)           & 14.13 & 13.58 & 0.31 \\
\textbf{GA (P=100)} & \textbf{13.92} & \textbf{13.60} & 0.61 \\
PSO                 & 17.08 & 14.07 & 0.20 \\
GWO                 & 15.48 & 14.72 & 0.43 \\
Hybrid              & 16.19 & 14.50 & 0.25 \\
\bottomrule
\end{tabular}
\end{table}

\section{Eredmények elemzése}

Az eredmények elemzésekor erős szezonális hatást figyelhetünk meg. Nyáron a magas napelem-termelésnek köszönhetően a napi költség drasztikusan, mintegy 67\%-kal csökkenthető (12.48-ról 3.93 RON-ra), míg télen a megtakarítás mértéke alacsonyabb, körülbelül 18\% (42.09-ről 34.42 RON-ra). A borús nap a kettő közötti átmenetet képezi. Az algoritmusok teljesítménye is eltérő: a GA (P=100) a komplex, változékony esetekben bizonyult a legrobusztusabbnak, míg a GWO télen, a szűkös erőforrások mellett adta a legjobb átlagot. A PSO gyorsasága ellenére nagy keresési térben hajlamos beragadni, a hibrid megoldás pedig kiegyensúlyozott teljesítményt nyújtott, bár a többletidő nem minden esetben eredményezett jobb optimumot.

\section{Implementációs részletek}

A hibrid megoldás PSO és GA fázisokra bontja az optimalizálást. A vezérlő logika egy részlete a \ref{code:hybrid_snippet}.\ kódrészletben látható.
\begin{lstlisting}[language=Python, caption={A hibrid algoritmus vezérlő logikája}, label=code:hybrid_snippet]
class Hybrid(AlgorithmBase):
    def solve(self):
        # 1. fazis: PSO (felfedezes)
        pso_steps = int(self.max_iter * 0.6)
        pso = PSO(self.func, self.dim, self.lb, self.ub,
                  self.pop_size, pso_steps)
        gb, gb_fit = pso.solve()

        # 2. fazis: GA (finomhangolas)
        ga = GA(self.func, self.dim, self.lb, self.ub,
                self.pop_size, self.max_iter - pso_steps)
        # ... tovabbi lepesek
\end{lstlisting}

\section{A 2×2-es benchmark protokoll}

A projekt jelenlegi állapotában a klasszikus háromszcenáriós futtatás mellett egy új, nap-indexelt benchmark pipeline is működik, amely valós inverteres PV profilt és valós ENTSO-E day-ahead árakat használ \cite{Ismail2014}. A benchmark fixen 16 reprezentatív napot vesz figyelembe, és minden napra két adatváltozatot épít: \texttt{real\_pv} (valós PV + generált load + valós/fallback ár) és \texttt{synthetic} (generált PV + generált load + valós/fallback ár). A teljes futásszám az előző fejezetben leírt $N_{eval} = 640$ algoritmusfuttatás.

\section{A 2×2-es kísérleti mátrix eredményei}

Az összegzett eredményeket a \ref{tab:realpv_synth_summary}.\ táblázat foglalja össze, amely jól mutatja az algoritmusok minőség--sebesség kompromisszumát a különböző zajszintek mellett.

\begin{table}[H]
\centering
\caption{A 2×2-es benchmark fő eredményei (640 futtatás aggregációja)}
\label{tab:realpv_synth_summary}
\resizebox{\textwidth}{!}{
\begin{tabular}{llccc}
\toprule
\textbf{Adathalmaz} & \textbf{Algoritmus} & \textbf{Átlagköltség (RON)} & \textbf{Gap \%} & \textbf{Győzelmi arány \%} \\
\midrule
\multirow{5}{*}{\textbf{(1) Lab-ideális}}
& \textbf{GA (P=100)} & \textbf{10.50} & \textbf{1.56} & \textbf{62.5} \\
& GA (P=50)  & 10.86 &  4.69 & 25.0 \\
& GWO        & 11.37 & 10.14 & 12.5 \\
& Hybrid     & 13.04 & 27.45 &  0.0 \\
& PSO        & 14.02 & 41.00 &  0.0 \\
\midrule
\multirow{5}{*}{\textbf{(2) Termelési zaj}}
& \textbf{GA (P=100)} & \textbf{13.61} &  \textbf{0.96} & \textbf{62.5} \\
& GA (P=50)  & 13.90 & 32.08 & 18.8 \\
& GWO        & 14.23 & 60.65 & 18.8 \\
& Hybrid     & 15.43 & 82.20 &  0.0 \\
& PSO        & 16.70 &159.08 &  0.0 \\
\midrule
\multirow{5}{*}{\textbf{(3) Fogyasztási zaj}}
& \textbf{GA (P=100)} &  \textbf{7.76} &  \textbf{2.73} & \textbf{62.5} \\
& GA (P=50)  &  7.96 &  3.72 & 31.2 \\
& GWO        &  8.77 & 13.15 &  6.2 \\
& Hybrid     &  9.79 & 23.30 &  0.0 \\
& PSO        & 11.14 & 45.33 &  0.0 \\
\midrule
\multirow{5}{*}{\textbf{(4) Valós világ}}
& \textbf{GA (P=100)} & \textbf{10.80} &  \textbf{1.97} & \textbf{50.0} \\
& GA (P=50)  & 11.13 &  5.24 & 31.2 \\
& GWO        & 11.64 & 17.29 & 18.8 \\
& Hybrid     & 13.17 & 44.28 &  0.0 \\
& PSO        & 13.92 & 72.32 &  0.0 \\
\bottomrule
\end{tabular}
}
\end{table}

A táblázatból egyértelműen látszik, hogy a GA (P=100) mind a négy adathalmazon robusztus maradt: az átlagos költsége és a gap\% értéke is konzisztensen a legjobb. A negyedik forgatókönyvben, ahol mind a termelés, mind a fogyasztás magas zajszinttel rendelkezett, a GA(P=100) győzelmi aránya 62.5\%-ról 50.0\%-ra csökkent, ami jelzi a keresési tér drasztikus nehezedését. A PSO a laboratóriumi körülmények között még elfogadható lemaradással futott, de a valós, zajos idősorok bevezetésekor a teljesítménye összeomlott -- a valós PV esetén akár 159\%-os gap-et is elért. Ez a benchmark megerősíti, hogy a valós adatok használata kritikus lépés a HEMS algoritmusok tesztelésében.

\section{A keresési tér és az algoritmusok viselkedése}

A 640 futásos benchmark egyik legszembetűnőbb eredménye, hogy a PSO szintetikus, laboratóriumi környezetben még versenyképes volt, azonban a valós fogyasztási adatok és a valós inverter-mérések együttes bevezetésekor a teljesítménye drasztikusan romlott. A generált, szintetikus profilok folytonos, sima görbéket eredményeznek, amelyek a célfüggvény terét egy viszonylag jól strukturált, közel konvex jellegű térbe rendezik. Ezzel szemben a valós UCI-alapú fogyasztási adatok a háztartási gépek véletlenszerű bekapcsolásából származó hirtelen tüskéket és aszimmetriát visznek a bemenetbe, ami a keresési teret rücskössé, több lokális optimumot tartalmazóvá alakítja.

A PSO-ban, ha a raj egy zajos idősor miatt egy hamis lokális optimumhoz vonzódik, a $\vec{g}_{best}$ komponens gyorsan magához rántja a többi részecskét, csökkentve a diverzitást és stagnálást okozva. A genetikus algoritmus ezzel szemben a mutációs operátor révén jobban fenntartja a populáció sokszínűségét: a véletlenszerű, gradiensfüggetlen ugrások segítenek abban, hogy az algoritmus ne ragadjon be tartósan egy lokális csapdába. Ebből az a mérnöki következtetés vonható le, hogy a HEMS-feladatoknál nem csak a gyors konvergencia, hanem a diverzitás fenntartása is kulcsfontosságú.

\section{Érvényességi korlátok}

A jelenlegi benchmark erős összehasonlító alapot ad, de néhány korlát tudatosan kezelendő. A futásszám (2 futás/profil) gyors iterációra optimalizált, nem végleges statisztikai mélységre, ezért a rangsorok megbízható mérnöki indikátorok, de végleges tudományos következtetésekhez érzékenységvizsgálat és bővebb adatkészlet szükséges. Az akkumulátor degradáció jelenleg nincs expliciten beárazva a célfüggvénybe, ami a gazdasági becslések pontosságát korlátozza.